\subsection{Automated AI-Firmware Integration Pipeline}

The framework uses a specialized integration node that resolves the directory structure of both the generated AI code and the target firmware project. Using shell-level primitives, the system synchronizes the source files (e.g., \texttt{network.c}, \texttt{network\_data.c}) into the \texttt{/Src} directory and the header files (\texttt{network.h}, \texttt{network\_data.h}) into the \texttt{/Inc} directory of the STM32 project. To achieve a truly end-to-end autonomous workflow, I implemented an automated integration pipeline.

\subsubsection{Source Code Synchronization}

The framework uses a specialized integration node that resolves the directory structure of both the generated AI code and the target firmware project. Using shell-level primitives, the system synchronizes the source files (e.g., \texttt{network.c}, \texttt{network\_data.c}) into the \texttt{/Src} directory and the header files (\texttt{network.h}, \texttt{network\_data.h}) into the \texttt{/Inc} directory of the STM32 project.

This synchronization ensures that the compiler can immediately resolve the new AI symbols without manual configuration of search paths in the IDE.

\subsubsection{LLM-Guided Code Injection in main.c}

The framework utilizes a combination of regex-based transformations and \texttt{sed} commands to insert the necessary include directives. For example, to inject the network headers, the system identifies the \texttt{/* USER CODE BEGIN Includes */} block and appends the appropriate headers as shown in Listing~\ref{lst:sed_injection}. My pipeline exploits these markers to inject code safely.

The framework utilizes a combination of regex-based transformations and \texttt{sed} commands to insert the necessary include directives. For example, to inject the network headers, the system identifies the \texttt{/* USER CODE BEGIN Includes */} block and appends the appropriate headers as shown in Listing~\ref{lst:sed_injection}.

\begin{lstlisting}[language=bash, caption={Automated code injection using sed.}, label={lst:sed_injection}, mathescape=false]
# Automated injection of AI headers into main.c
sed -i '' '/\/\* USER CODE BEGIN Includes \*\//a\
#include "network.h"\
#include "network_data.h"
' "$PROJ_DIR/Src/main.c"
\end{lstlisting}

Furthermore, the higher-level agentic logic analyzes the \texttt{main.c} structure to determine the optimal insertion point for the model initialization call (e.g., \texttt{ai\_network\_create} and \texttt{ai\_network\_init}). By automating these steps, the integration time is reduced from several minutes of manual error-prone work to a sub-second atomic operation, enabling the rapid validation of dozens of model variants within the same firmware baseline.
